\paragraph{Short-horizon prediction vs.\ trip-level route generation.}
Most trajectory prediction work targets short horizons and evaluates point-wise accuracy, where committing to distinct trip-scale corridors is not required.
In contrast, route generation must represent corridor-level alternatives under the same OD intent, while remaining feasible on a transportation network.

\paragraph{Foundation models and transfer for mobility.}
Recent trajectory foundation models \citep{unitraj2025,lin2024trajfm} aim to learn reusable representations from large-scale traces.
These approaches are complementary: we do not treat \emph{route choice} as a reconstruction problem, but as a context-conditioned decision whose distribution should shift with time and urban semantics even for the same OD pair.
Our work therefore focuses on decision mechanisms and distributional evaluation at the corridor level.

\paragraph{Map-conditioned and road-network constrained generation.}
Multiple lines of work incorporate road structure to improve physical fidelity, either by conditioning on street maps \citep{ijcai2025p1041} or by constraining diffusion models on road networks \citep{wei2024diffrntraj,zhu2024controltraj,yang2025geo}.
Our diagnostic evidence aligns with this direction: map-free continuous generators can exhibit off-road artifacts at trip scale.
We go one step further by making corridor choice explicit as a discrete decision and by treating temporal--spatial semantics as first-class conditioning signals for \emph{behavioral} realism.

\paragraph{Coarse-to-fine and path-generation diffusion.}
Coarse-to-fine generation has emerged as a practical strategy for long-horizon mobility modeling \citep{guo2025cardiff}, and latent diffusion has also been explored for road-network paths \citep{fan2024diffpath}.
CascadeTraj follows this engineering intuition but targets a different unit of decision: rather than denoising coordinates toward a measure-zero road manifold, we commit to sparse graph-level decisions and let constrained execution guarantee connectivity.

\paragraph{Cross-city generalization.}
GTG \citep{wang2025gtg} demonstrates that trajectory generation can generalize across cities.
We leverage the same multi-city setting but ask a more specific question: what modulates corridor preference under a shared road topology?
This motivates our emphasis on explicit context (time, road hierarchy, and urban semantics) and corridor-distribution evaluation beyond shortest-path biases.
